% Geraden · Punktprobe und Punkte auf der Geraden · Prüfungs-Fokus Abitur GK – bank/geraden/e2.jsonl (zusammenbau v0.6, Prüfungs-Fokus)
\einheitenkopf[e2]{Einheit 2 · Punktprobe und Punkte auf der Geraden}
\zweigzeile{Abi GK ×7}
\setcounter{aufgabe}{0}

% Anlauf (ohne Prüfkennung)
\begin{aufgabe}{Punktprobe und Punkte auf der Geraden – Anlauf}
\begin{teile}
\teil Für $g\colon \vec x = (2 | 5 | 1) + t \cdot (0 | 3 | 4)$ und $P(2 | 11 | 9)$: Aus welcher Koordinate lässt sich $t$ bestimmen? \kreuz{nur aus der ersten} \\ \kreuz{aus der zweiten oder der dritten} \\ \kreuz{aus keiner}
\teil Prüfe, ob $P(7 | 5 | 8)$ auf $g\colon \vec x = (1 | 2 | -1) + r \cdot (2 | 1 | 3)$ liegt.
\end{teile}
\end{aufgabe}

\begin{aufgabe}{Punktprobe und Punkte auf der Geraden – Prüfungsaufgaben}
\begin{teile}
\teil Gegeben sind $g\colon \vec x = (5 | -1 | 2) + t \cdot (3 | 0 | -2)$, $t \in \mathbb{R}$, und $A(2 | 3 | 4)$. Begründe, dass $A$ nicht auf $g$ liegt. \hfill (A23)
\teil Gegeben sind $h\colon \vec x = (0 | 4 | 6) + r \cdot (2 | -3 | 0)$, $r \in \mathbb{R}$, und $B(4 | -2 | 5)$. Begründe, dass $B$ nicht auf $h$ liegt. \hfill (A23)
\teil Gegeben sind $k\colon \vec x = (-3 | 2 | 1) + t \cdot (0 | 1 | 4)$, $t \in \mathbb{R}$, und $C(1 | 5 | 13)$. Begründe, dass $C$ nicht auf $k$ liegt. \hfill (A23)
\teil Der Punkt $B(-2 | y_B | z_B)$ liegt auf $g\colon \vec x = (1 | 3 | 2) + r \cdot (-1 | 2 | -3)$. Bestimme $y_B$ und $z_B$. \hfill (A19) \\ $y_B$ = \leerfeld , $z_B$ = \leerfeld
\teil Der Punkt $P(x_P | 2 | z_P)$ liegt auf $g\colon \vec x = (0 | -1 | 4) + t \cdot (2 | 1 | -2)$. Bestimme $x_P$ und $z_P$. \hfill (A19) \\ $x_P$ = \leerfeld , $z_P$ = \leerfeld
\teil Der Punkt $Q(x_Q | y_Q | -3)$ liegt auf $h\colon \vec x = (3 | -2 | 0) + t \cdot (-2 | 4 | 1)$. Bestimme $x_Q$ und $y_Q$. \hfill (A19) \\ $x_Q$ = \leerfeld , $y_Q$ = \leerfeld
\teil Die Gerade $g$ verläuft durch $A(3 | 1 | 2)$ mit dem Richtungsvektor $\vec u = (2 | -1 | 2)$. Bestimme die Koordinaten eines Punkts auf $g$, der von $A$ den Abstand 9 hat. \hfill (A19)
\teil Die Gerade $g$ verläuft durch $A(0 | 4 | -1)$ mit dem Richtungsvektor $\vec u = (4 | 0 | -3)$. Bestimme die Koordinaten eines Punkts auf $g$, der von $A$ den Abstand 10 hat. \hfill (A19)
\teil Die Gerade $g$ verläuft durch $A(-2 | 1 | 3)$ mit dem Richtungsvektor $\vec u = (1 | 2 | 2)$. Bestimme die Koordinaten eines Punkts auf $g$, der von $A$ den Abstand $4{,}5$ hat. \hfill (A19)
\teil Eine Hecke verläuft geradlinig von $E(8 | 0 | 0)$ nach $F(0 | 8 | 0)$ (1 LE = 1 m). Weise nach, dass der Schattenpunkt $T(2 | 6 | 0)$ auf der Strecke $EF$ liegt. \hfill (A18)
\teil Ein Zaun verläuft geradlinig von $E(6 | -2 | 0)$ nach $F(2 | 2 | 0)$ (1 LE = 1 m). Entscheide, ob der Schattenpunkt $T(-3 | 7 | 0)$ auf dem Zaun $EF$ liegt. \hfill (A18)
\teil Ein Stahlseil ist geradlinig von $A(1 | 1 | 4)$ nach $B(7 | 4 | 1)$ gespannt (1 LE = 1 m). Weise nach, dass der Punkt $L(5 | 3 | 2)$, an dem eine Lampe hängt, auf dem Seil liegt. \hfill (A18)
\end{teile}
\end{aufgabe}

\begin{aufgabe}{Punktprobe und Punkte auf der Geraden – Prüfungsaufgaben – weiter}
\begin{teile}
\teil Weise nach, dass die Punkte $A(2 | -1 | 3)$, $B(4 | 0 | 1)$ und $C(8 | 2 | -3)$ auf einer Geraden liegen. \hfill (A20)
\teil Weise nach, dass die Punkte $A(0 | 5 | 2)$, $B(1 | 3 | 3)$ und $C(-2 | 9 | 0)$ auf einer Geraden liegen. \hfill (A20)
\teil Untersuche, ob die Punkte $A(1 | 0 | 2)$, $B(3 | 1 | 5)$ und $C(7 | 3 | 12)$ auf einer Geraden liegen. \hfill (A20)
\teil Ein Sessellift führt geradlinig von $A(1 | 20 | 2)$ nach $E(7 | 2 | 11)$; die Talstation liegt im Ursprung, die $x_1x_2$-Ebene ist waagerecht, 1 LE entspricht 5 m. Wie hoch über der Talstation ist der Lift an der Stelle, die 84 m vom Anfang $A$ entfernt ist? \hfill (A23) \\ \leerfeld[m]
\teil Eine Materialseilbahn führt geradlinig von $A(0 | 0 | 4)$ nach $E(8 | 4 | 12)$; der Ursprung liegt auf Höhe der Talstation, die $x_1x_2$-Ebene ist waagerecht, 1 LE entspricht 10 m. Wie hoch über der Talstation ist die Bahn 30 m nach dem Anfang $A$? \hfill (A23) \\ \leerfeld[m]
\teil Eine Standseilbahn fährt geradlinig von $A(3 | 1 | 0)$ nach $E(-1 | 9 | 8)$; die Talstation liegt im Ursprung, die $x_1x_2$-Ebene ist waagerecht, 1 LE entspricht 20 m. Wie hoch über der Talstation ist die Bahn 180 m nach dem Anfang $A$? \hfill (A23) \\ \leerfeld[m]
\teil Ein Prisma hat die Ecken $A(5 | 0 | 0)$ und $G(1 | 4 | 6)$. Gib eine Gleichung der Geraden durch $A$ und $G$ an und weise nach, dass der Punkt $S(3 | 2 | 5)$ nicht auf ihr liegt. \hfill (A26)
\teil Ein Quader hat die Ecken $A(0 | 6 | 1)$ und $G(4 | 0 | 4)$. Gib eine Gleichung der Geraden durch $A$ und $G$ an und weise nach, dass der Punkt $S(2 | 3 | 2)$ nicht auf ihr liegt. \hfill (A26)
\teil Ein Prisma hat die Ecken $A(2 | 0 | 0)$ und $G(-2 | 6 | 6)$. Gib eine Gleichung der Geraden durch $A$ und $G$ an und untersuche, ob der Punkt $S(0 | 3 | 3)$ auf ihr liegt. \hfill (A26)
\end{teile}
\end{aufgabe}

\medskip
%% Merkkasten Einheit 2: 7 Zeilen, Vorgabe höchstens fünf (3.1) – ungekürzt
\uebersichtskasten{Punktprobe: den Punkt mit der Geradengleichung gleichsetzen, den Parameter aus einer Koordinate bestimmen und in den übrigen prüfen – ein Widerspruch heißt: der Punkt liegt nicht auf g; kein Widerspruch: er liegt darauf. Bei einer Strecke zusätzlich prüfen, ob der Parameter zwischen null und eins liegt. \\ g: x = (8 | 3 | −3) + s · (−4 | 0 | 3), P(4 | 3 | 3): aus der ersten Koordinate s = 1, aus der dritten s = 2 – Widerspruch, P liegt nicht auf g. \\ Schnell entschieden: ist eine Richtungskomponente null, hat jeder Geradenpunkt dort dieselbe Koordinate – weicht der Punkt ab, ist die Probe ohne Rechnung entschieden. \\ Jeder Punkt von g hat die zweite Koordinate 3. \\ Punkt mit Bedingung: den Parameter aus der Bedingungskoordinate bestimmen und einsetzen. \\ Punkt mit Abstand: der Abstand d vom Aufpunkt liegt bei p ± d/|u| · u – erst den Richtungsvektor normieren, dann strecken. \\ Drei Punkte auf einer Geraden: Gerade durch zwei aufstellen und den dritten per Probe prüfen – oder die Verbindungsvektoren auf Kollinearität.}
